\section{Related Work}
\label{sec:related}

\textbf{Surveys of patient pathway mining.} There exist two recent surveys that provide a good overview of this area. Rosa et al. \cite{rosa2022modelling} provide a survey of the models and techniques used for patient pathways mining. They distinguish between the sequence mining and process mining approaches. Flothow et al. \cite{flothow2023} provide a survey of various techniques applied in the search of patient sequences of healthcare use, distinguishing clustering approaches, which uncover entire patient trajectories, from pattern mining approaches, which are applied to subsequences.

\textbf{Combining clustering with sequence mining.} The most relevant research hybridizes the two concepts. Najjar et al. \cite{najjar2018} suggest a two-step process, whereby they cluster single healthcare visits using K-Prototypes and derive treatment pathways afterwards. It is an important predecessor of our research because it also uses the mixed clustering technique that we employ. Baudrier et al. \cite{baudrier2022} hybridize their care trajectory clustering on patients with kidney cancer with a sequential pattern discovery algorithm and conclude that their clusters differ from each other by actual results like mortality rate. Vogt et al. \cite{vogt2018} discover care pathway clusters for heart-failure patients from sequences in insurance claims with ambulatory healthcare pathways. Our approach is conceptually similar, yet it differs from theirs in the order of steps involved.

\textbf{Sequence mining from EMR with interpretation.} Perer et al. \cite{perer2015} conduct sequence mining from electronic medical records based on frequent event sequences to patient outcomes, with the major emphasis on visual analysis. Their results show that frequent sequences can detect clinically meaningful care patterns unknown to the physician, which motivates our own attention to pattern quality, since frequency alone is not enough.

\textbf{Clustering of patients on EHR data.} Some studies concentrate on analyzing methods of proper patients' clustering. Yang et al. \cite{yang2024} develop a smart-healthcare solution aimed at clustering patients based on different kinds of information contained in their medical records, to reveal potential health risks and disease patterns across patient groups. Aljohani \cite{aljohani2024} evaluates eight different clustering methods for patients' stratification based on EHR and analyzes them using a multi-criteria approach.

\textbf{Process mining as an alternative.} An alternative set of methods analyzes care processes from event logs. Specifically, Murazzano et al. \cite{murazzano2026} use process mining to analyze chronic respiratory patients and observe a large variability in their care processes, pointing to another direction of analysis that addresses the same objective in a different way.

Clustering and sequence mining are both mature domains in healthcare analytics. There have been some research works which try to link them together. However, there remains the need for a unified framework that integrates cluster-scope mining with a rank of clinically meaningful patterns followed by an AI-based explanation component. That's where HospitAI comes in.